\section{引言}\label{sec:introduction}

一个置换约束规定：对另一端点的每个标记，本端点恰有一个对应标记。
唯一游戏猜想断言：即使几乎所有约束都能被满足，这类约束系统也可能
难以近似。我们正面解决该猜想，给出一个显式的无权归约，从
\(\mathrm{3SAT}\) 出发。

\subsection{定理}

\emph{唯一游戏实例}由有限顶点集、有限字母表 \(K\)、以及非空的定向
约束列表 \(e=(u_e,v_e,\pi_e)\) 组成，其中 \(\pi_e\) 是 \(K\) 的一个
置换。染色 \(a\) 满足 \(e\) 当
\(a(v_e)=\pi_e(a(u_e))\)。其值是被满足约束的比例，
\(\val(G)\) 是染色的最大值。当顶点与约束列表被完整写出（包括每个
置换的表）时，称实例\emph{显式且无权}。简单二部实例有两个不相交的
顶点集，且每对相对顶点上至多一条约束。

\begin{theorem}[唯一游戏定理]\label{thm:ugc}
对每个固定的 \(\varepsilon,\delta\in(0,1/2)\)，存在整数 \(s\ge1\)
与一个确定性多项式时间归约，从 \(\mathrm{3SAT}\) 到 \(K=\F_2^s\) 上
显式无权唯一游戏实例 \(G_\varphi\)，使得
\[
\begin{aligned}
 \varphi\text{ 可满足}
 &\quad\Longrightarrow\quad \val(G_\varphi)\ge1-\varepsilon,\\
 \varphi\text{ 不可满足}
 &\quad\Longrightarrow\quad \val(G_\varphi)\le\delta.
\end{aligned}
\]
图为简单二部图，且每条约束都是 \(K\) 的一个平移
\(a(v_e)=a(u_e)+c_e\)。字母表、运行时间多项式的次数及其常数只依赖
\(\varepsilon,\delta\)。
\end{theorem}

因此，该结果以其通常的边值表述正面解决唯一游戏
猜想~\cite[Conjecture~2.5]{Khot10}。两个误差被独立指定，且先于
字母表选取。二元平移约束是该猜想的一个已知的等价表述
~\cite[Corollary~13 and Section~11.1]{KKMO07}；本归约直接产生该
形式。

Theorem~\ref{thm:ugc} 为已有的针对割、覆盖、约束满足、排序、删除与
聚类问题的难度归约提供了唯一游戏假设
\cite{KKMO07,KhotRegev2008,Raghavendra08,GuruswamiEtAl2011OrderingCSP,CKKRS2006,DemaineEtAl2006Clustering}。
那些归约及由此得到的近似阈值均出自上述作者。
Section~\ref{sec:consequences} 给出精确的问题表述、严格阈值与归约
限定条件，并附各推论的证明。

\subsection{证明、编码与量词}

核心难点在于：当每个标记只容许唯一应答时，如何保持接近 1 的完备
性。证明验证与近似难度之间的早期联系由 Feige、Goldwasser、
Lov\'asz、Safra 与 Szegedy~\cite{FGLSS96}，由 Arora 与 Safra 经由
证明复合~\cite{AroraSafra98}，以及由 Arora、Lund、Motwani、Sudan
与 Szegedy 经由常数查询 PCP 定理~\cite{ALMSS98} 发展而来。
Bellare、Goldreich 与 Sudan 为验证者复合引入了长码（long code）与
折叠~\cite[Section~3.3]{BGS98}。答案 \(a\) 的长码是表
\(f\mapsto f(a)\)，以答案集上的布尔函数为指标。在 H\aa stad 的奇偶
难度证明中，傅里叶分析从使验证者足够频繁接受的任意证明表恢复出
外层博弈策略，而无需先恢复邻近码字。他的近可满足奇偶间隙是我们的
起始难度输入~\cite{Hastad01}。

Khot 于 2002 年提出唯一游戏猜想，并给出基于半定舍入的算法
~\cite{Khot02}。若干算法保证可以解释其量词顺序为何重要。对固定的
目标误差，Trevisan 把半定舍入与图分解结合，在输入误差按实例规模的
倒数对数速率缩小时得到多项式时间保证~\cite{Trevisan08}。
Charikar、Makarychev 与 Makarychev 用随机化半定舍入，从值至少
\(1-\epsilon\) 的 \(q\) 标记实例得到期望值
\(1-O(\sqrt{\epsilon\log q})\)~\cite{CMM06}；该保证随字母表增大而
失准，而猜想允许在误差选定后使用大的固定字母表。Kolla 的谱枚举
给出以标记扩展图（顶点为「变量–标记」对）适当谱子空间维数计的运行
时间界~\cite{Kolla11}。Arora、Barak 与 Steurer 用谱分解与枚举，对
足够小的固定完备性误差得到近似满足解，其保证与固定字母表规模无关
~\cite{ABS10}；对固定的正误差，其运行时间界是次指数的而非多项式
的。这些是已知保证的各自局限。

整性间隙关心的是松弛的强度。Khot 与 Vishnoi 构造了半定值任意接近
1、整数值任意接近 0 的唯一游戏实例，即使加上三角不等式也如此
~\cite{KV15}。Barak、Gopalan、H\aa stad、Meka、Raghavendra 与
Steurer 基于 Reed--Muller 码引入短码（short code），作为长码在字母
表归约与定量整性间隙等应用中的更短替代~\cite{BGHMRS15}。这些构造
约束的是特定松弛；难度归约需要单独论证。Bafna、Barak、Kothari、
Schramm 与 Steurer 对具有适当低度平方和超压缩性证书的图（含噪声
超立方与噪声 Reed--Muller 短码图）的近可满足平移博弈，得到多项式
时间常值保证~\cite{BBKSS21}。下文使用的秩一矩阵短码是另一种图，且
我们的论证不假设这些证书。

\subsection{Grassmann 扩张与完备性障碍}

最接近的前置难度突破涉及完备性不完美的二对一与二对二博弈。二对一
约束是每个右标记恰有两个原像的投影。二对二约束把两侧标记各分成
对，并在这些对之间建立双射匹配，允许匹配对内的四种组合。Khot、
Minzer 与 Safra 的 STOC 2017 工作提出 Grassmann 图归约，其非标准的
可靠性保证以一个结构性假设为条件~\cite{KMS17}。Dinur、Khot、
Kindler、Minzer 与 Safra 的 STOC 2018 工作提出具有普通边满足可靠性
的归约，以一致性假设为条件~\cite{DKKMS25}；其扩张工作提出了相关的
结构问题并建立了部分情形~\cite{DKKMS21Expansion}。Barak、Kothari
与 Steurer 通过二元矩阵表示把扩张与所需的一致性测试联系起来
~\cite{BKS18}。Khot、Minzer 与 Safra 的扩张定理于 2018 年完成了
这条线~\cite{KMS23}。

这些结果对几乎可满足的二对一与二对二博弈建立了难度，并作为推论给
出完备性接近二分之一、可靠性任意小的唯一游戏间隙。使唯一游戏的
完备性任意接近 1 是一个单独的问题。\footnote{一篇独立的姊妹工作
证明：在固定字母表上，完备性完美、固定小可靠性下精确二对一博弈的
难度~\cite[Theorem~1.1]{OpenAIPerfectTwoToOne2026}。它不是本证明的
输入。}

Barak、Kothari 与 Steurer 的二次矩阵短码~\cite{BKS18} 是本文使用
的表述：其顶点是二元矩阵，其步骤添加秩一矩阵。他们的双应答规则
容许旧答案或其被更新左因子平移后的答案；只要求相等只能给出二分之
一的完备性保证。他们讨论过更高次的测试，作为改进完备性的可能路线
~\cite[Section~1.6]{BKS18}。我们使用的逆命题（由 Grassmann 扩张
导出）见 Theorem~\ref{thm:shortcode}。Ellis、Kindler 与 Lifshitz
随后给出更短的扩张证明，使用线性映射空间上的超压缩性~\cite{EKL26}；
我们的外部输入仍是 Khot--Minzer--Safra 定理。我们的证明保留已有的
二次逆定理，并用一个非线性映射改变噪声。下面的比较解释这一选择。

\subsection{构造与证明}

均匀秩一测试展示了完备性问题。设 \(M\) 是二元 \(\ell\times m\)
矩阵，固定非零答案 \(z\in\F_2^m\)。求值函数 \(f_z(M)=Mz\) 满足
\[
 f_z(M+a l^\top)-f_z(M)=a(l^\top z).
\]
对独立的均匀 \(a\in\F_2^\ell\) 与 \(l\in\F_2^m\)，该求值以概率
\((1+2^{-\ell})/2\) 保持其值，接近二分之一。我们的测试改变噪声，
同时保留普通秩一测试作为可靠性分析的结构工具。

对任意指定的 \(p>0\)，Lemma~\ref{lem:latent} 固定一个秩阈值，并对
每个足够大的 \(\ell\) 构造包含 \(K=\F_2^\ell\) 的二元空间
\(\mathcal V\)、映射 \(C:\mathcal V\to K\) 与 \(\mathcal V\) 上的
噪声律 \(\mu\)，使得对 \(a\sim\mu\)，
\[
 C(x+h)=C(x)+h\quad(h\in K),\qquad
 \Pr[C(x+a)\ne C(x)]\le p
\]
其中 \(x\) 均匀且与 \(a\) 独立。同时，任何限制到 \(K\) 后秩足够大
的线性观察族以至少 \(1/8\) 的概率检测出 \(a\)。秩阈值只依赖
\(p\)，且在 \(\ell\) 之前选定。该构造使用有限域上的二次块：其非
线性输出在递归下越来越稳定，而一个秩势控制线性观察的损失。
Section~\ref{sec:latent} 证明这一分离，并把输出平移实现为子空间
\(K\)。

对同一个非零 \(z\)，设 \(P\in\Hom(\F_2^m,\mathcal V)\) 均匀。非线性
求值 \(P\mapsto C(Pz)\) 在更新 \(P\mapsto P+a l^\top\) 下比较
\[
 C(Pz)\quad\text{与}\quad C\bigl(Pz+a(l^\top z)\bigr)
\]
两个值。此处 \(Pz\) 在 \(\mathcal V\) 中均匀，与潜在噪声 \(a\)
独立，而 \(l^\top z\) 是独立的公平比特。因此其相等测试以至多
\(p/2\) 的概率失败。

归约从三个不同变量上的二元奇偶方程出发。每条方程有四个满足的
三元组。方程元组 \(U=(e_1,\ldots,e_k)\) 上的表是从其解平面的乘积
到 \(\mathcal V\) 的仿射映射，由一个 \(m=1+2k\) 列的矩阵表示：一个
截距列加每条方程两个斜率列。当某个方程块只通过一个实际比特影响
一张表时，其约化键（reduced key）记录该变量与整个诱导仿射映射。
投影表与其拉回具有相同的键。把键的常数 \(K\)-平移归组为一个顶点，
就折叠了答案：代表元处的标记决定每个平移后的答案，而共享的键在
两个域上给出相同的恢复答案。所得的唯一游戏测试利用潜在噪声，检验
一张表与其秩一扰动上恢复答案的相等性。一个近满足的全局奇偶赋值在
其满足的每个元组上给出上述理想求值，从而提供接近 1 的完备性。
键共享改编了 Khot 与 Safra 的虚拟表~\cite[Section~1.1]{KS13}，而
平移轨道改编了 Bellare、Goldreich 与 Sudan 的折叠原理
~\cite[Section~3.3]{BGS98}。

对可靠性，设某染色以至少 \(0.99\) 的概率通过该测试。矩阵频率是
元组 \(g=(g_i)_{i=0}^{m-1}\in(\mathcal V^*)^m\)（每列一个泛函），
其特征为 \((-1)^{\sum_i g_i(P\mathbf e_i)}\)，其中 \(\mathbf e_i\)
是 \(\F_2^m\) 的标准基向量。把秩一更新的相位增量对均匀选择子
\(l\) 取平均：
\[
\begin{aligned}
 \E_{a\sim\mu,l}(-1)^{\sum_i l_i g_i(a)}
   &=\Pr_{a\sim\mu}[g_i(a)=0\text{ 对每个 }i],\\
 \E_{a\text{ 在 }K\text{ 中均匀},l}(-1)^{\sum_i l_i g_i(a)}
   &=2^{-\dim\Span\{g_i|_K\}}.
\end{aligned}
\]
第一个乘子非负，且由潜在检测，当限制秩大时至多 \(7/8\)。因此高
接受率迫使固定数量的低秩傅里叶质量；第二个乘子在 \(K\) 中的普通
均匀秩一扰动下把这一质量转化为正的接受率。使用与潜在数据一起固定
的线性补 \(\mathcal V=K\oplus W\)，把表写成 \(P=(M,T)\)，其中
\(M\) 与 \(T\) 是其取值于 \(K\) 与 \(W\) 的仿射分量，用含常数列的
矩阵表示。普通扰动改变 \(M\) 而固定 \(T\)。于是，对正比例的对
\((U,T)\)，作为 \(M\) 的函数的恢复答案经常在普通测试下保持取值。
对其余的可靠性论证，取 \(k\) 为足够大的整数立方，潜在与短码参数
固定。在每个这样的对上，逆短码定理给出与某个目标
\(Mz+u\)（\(z\in\F_2^m\)、\(u\in K\)）的一致性，该一致性在由至多
固定数目 \(r_{\rm s}\) 个行方程与至多 \(r_{\rm s}\) 个列方程定义的
切片上成立。

对独立的均匀线性映射 \(A:K\to\F_2^{r_{\rm s}}\)，关于 \((M,T)\)
的可见建议是 \((A,T,AM)\)。连同 \(U\)，它确定了 \(T\) 以及 \(U\)
的答案空间上 \(K\) 值仿射映射的行纤维 \(\{M':AM'=AM\}\)（含其常数
项）。一个擦除论证证明：在 \(U\) 上的无限制表采样下，固定正比例的
此类数据容许额外的列方程，以及一个在非空切片上见证一致性的目标
\(Mz+u\)。对每个这样的数据 \((U,A,T,AM)\)，按只依赖该可见数据与
固定染色的固定规则选取一个证书。折叠随后使该证书确定 \(U\) 上的
一个短的、非空的合法答案集。

建议实验把方程元组 \(U\) 与问题 \(O\) 比较：\(O\) 以概率 \(\beta\)
独立地把每条方程替换为一个均匀选取的变量。设 \(\pi\) 为实际答案的
所得投影。裁判在 \(O\) 上采样均匀表 \((M_O,T_O)\) 及具有上述律的
独立映射 \(A\)。在实际答案空间上，它把可见数据
\((A,T_O\circ\pi,AM_O\circ\pi)\) 连同 \(U\)、把 \((A,T_O,AM_O)\)
连同 \(O\) 发给两个玩家。取下文的稀疏 \(\beta\)，对投影选择的
取平均后，完整的拉回表律接近抽取所用的无限制表律。在正比例的联合
样本上，所选的 \(U\)-证书转移到 \(O\) 上的非空切片并保持一致性。

两个策略由各自的可见数据定义。当其建议容许证书时，\(U\)-策略从
其短答案集中均匀采样；在其他输入上返回任意合法答案。\(O\)-策略
在自己的行纤维 \(\{M':AM'=AM_O\}\) 上使用傅里叶解码器。在分析中，
转移来的证书提供傅里叶质量，证明该解码器以正概率到达投影后的
\(U\)-答案集；解码器本身不接收该证书。独立的 \(U\)-样本随后给出
完整答案一致性，概率至少为固定的 \(\gamma>0\)，与 \(k\) 无关。
稀疏问题比较与傅里叶解码改编了 Khot 与 Safra 的方法~\cite{KS13}；
相关建议与消失坐标的使用遵循 Khot、Minzer 与 Safra
~\cite[Sections~3.3--3.4]{KMS17}，以及 Minzer 与 Zheng 对该框架的
后续使用~\cite[Section~3]{MZ26}。

在 NO 奇偶实例上，每对建议策略的一致性趋于零。取 \(\beta=k^{-2/3}\)
使 \(k\beta^2\to0\)（表律比较所需），同时 \(k\beta\to\infty\) 留下
大量单例坐标。当每条被观察的行在其保留比特上的系数为零时，称该
单例是干净的。在分析中，以 \(A\)、干净位置、\(O\) 上的原建议截距、
以及其他位置上采样的问题与建议斜率为条件。干净问题对于是保持独立
的方程–变量律。对这些固定数据的每个取值，由每个原局部输入自身的
干净问题重构它，即定义了重复外层博弈的策略。并行重复界定它们在
干净坐标上的成功率；对那些坐标的个数取平均，把原建议一致性界定在
\(\exp(-c k^{1/3})\) 以内。对大 \(k\)，这与解码后的下界矛盾。

并行重复我们使用 Dinur 与 Steurer 的显式投影博弈界
~\cite[Corollary~1.2]{DS14}。Raz 对一般双证明者博弈建立了指数衰减
~\cite{Raz98}，Holenstein 简化了证明~\cite{Holenstein2009}；那些界
依赖答案字母表。Rao 对投影博弈去除了这一依赖~\cite[Theorem~4]{Rao11}。
Dinur--Steurer 界先用于上述干净坐标，最后用于把矩阵博弈权重取整并
把每条边细分为四条边之后。其字母表无关性允许最终重复指数在潜在
字母表之前固定，从而完成显式无权归约。

Figure~\ref{fig:overview} 展示该归约以及确立可靠性的两个不相容
界。
\begin{figure}[htbp]
\centering
\begin{tikzpicture}[
  x=1cm,y=1cm,
  every node/.style={font=\small},
  box/.style={draw=linkblue!80!black,rounded corners=2pt,
    text width=5.9cm,align=center,inner sep=8pt,minimum height=11mm},
  branch/.style={box,text width=2.70cm,inner sep=6pt,minimum height=15mm},
  heading/.style={font=\small\bfseries,align=center},
  reduction/.style={-{Latex[length=2.1mm]},thick,draw=linkblue!80!black},
  implication/.style={-{Latex[length=2.1mm]},thick,dashed,draw=black!70}
]
\node[heading] at (0,0) {归约};
\node[heading] at (7.4,0) {NO 输入上的可靠性};

\node[box] (parity) at (0,-1.0)
  {$\mathrm{3SAT}$\\$\longrightarrow$ 带权奇偶间隙实例};
\node[box] (matrix) at (0,-2.8)
  {潜在矩阵测试\\
   稳定非线性映射：Lemma~\ref{lem:latent}};
\node[box] (ugc) at (0,-6.2)
  {无权简单二部\\唯一游戏\\
   Theorem~\ref{thm:ugc}};
\draw[reduction] (parity) -- (matrix);
\draw[reduction] (matrix) -- node[midway,fill=white,inner sep=4pt,
    text width=5.0cm,align=center]
  {取整与细分；\\并行重复} (ugc);

\node[branch] (assumption) at (5.65,-1.0)
  {假设测试接受率\\至少 $0.99$};
\node[branch] (no) at (9.15,-1.0)
  {NO 奇偶\\实例};
\node[branch] (decoding) at (5.65,-3.3)
  {解码后的建议策略\\一致概率至少 $\gamma>0$\\
   Proposition~\ref{prop:matrix-decoding}};
\node[branch] (clean) at (9.15,-3.3)
  {每一对策略：\\一致概率至多\\$\exp(-c k^{1/3})$\\
   Lemma~\ref{lem:clean}};
\node[box] (contradiction) at (7.4,-6.2)
  {对足够大的 $k$ 得到矛盾：\\
   $\exp(-c k^{1/3})<\gamma$};
\draw[implication] (assumption) -- (decoding);
\draw[implication] (no) -- (clean);
\draw[implication] (decoding) -- (contradiction);
\draw[implication] (clean) -- (contradiction);
\end{tikzpicture}
\caption{实线箭头描述归约；虚线箭头描述其可靠性论证。解码下界来自
高测试接受率，而干净坐标上界来自 NO 奇偶间隙，并适用于每个建议
策略。常数 $c,\gamma>0$ 在元组长度 $k$ 增长之前固定。}
\label{fig:overview}
\end{figure}


Section~\ref{sec:preliminaries} 固定博弈与概率约定。
Section~\ref{sec:latent} 构造 gadget，Section~\ref{sec:game} 定义
置换测试。Sections~\ref{sec:matrix} 与~\ref{sec:clean} 给出建议
实验的两个不相容界。Section~\ref{sec:finish} 完成参数选取与显式
唯一游戏归约。Section~\ref{sec:consequences} 从完成的定理与所引
归约导出近似推论。Appendix~\ref{app:finite-consequences} 给出那里
使用的简单强左表述与有限排序归约。
